\scanpage{146}
Question : \unsure{Dans quel cas (uniforme convexité} \unsure{les|ces|des} \unsure{espaces|esp.} $L^p(E,F)$ \unsure{?} pour $0<p<\infty$ ? Pas d'importance.

\unsure{2|1}) Pour la continuité de $h\to h^r$, \unsure{on} est ramené au cas $0\leq r<1$ (car c'\unsure{est} \unsure{trivial} pour $r$ entier) Il semble possible que

\[
\bigl\|\lvert {\unsure{h|H}}^r-k^r\rvert\bigr\|^{1/r}\leq\|h-k\|\qquad\text{pour }0\leq r\leq1,
\]

Il revient au $=$ de \unsure{prouver|présumer} que

\[
S_{p/r}(h^r-k^r)=S_p\bigl(\lvert {\unsure{h|H}}^r-k^r\rvert^{1/r}\bigr)
\]

est fonction croissante de $r$, ($h$ et $k$ hermitiens positifs, $p$ quelconque)\unsure{.|,}

On peut faire aussi $p=1$.
